\documentclass[12pt,a4paper,twoside]{article}
\usepackage{geometry}
\geometry{a4paper}

\usepackage[utf8]{inputenc}
\usepackage[OT1]{fontenc}
\usepackage{polski}
\usepackage{graphicx}
\usepackage{setspace}
\usepackage{fancyhdr}
\usepackage{amsmath}
\usepackage{amssymb}

\begin{document}
\newgeometry{margin=3cm}
\onehalfspacing

\begin{titlepage}
	\begin{center}
	{\LARGE Uniwersytet im. Adama Mickiewicza w Poznaniu \\
	Wydział Matematyki i Informatyki }
	\line(1,0){350}

	\vspace{1cm}
	\includegraphics[width=2cm]{logo-uam/logo-uam.png}
	\vspace{1cm}

	\vspace{1cm}
	{\Huge Ataki na kryptograficzne \\ funkcje skrótu} \\[0.5cm]
	{\Large Marcin Kurczewski}
	\end{center}

	\vspace{3cm}
	\hspace{8cm}\parbox[l]{6cm}{\Large Praca magisterska \\
	napisana pod kierunkiem \\
	dr Michała Rena}

	\begin{center}
	\vspace{4cm}
	Poznań, maj 2013
	\end{center}

\end{titlepage}


\newpage
\thispagestyle{empty}
\begin{center}
	OŚWIADCZENIE
\end{center}

Ja, niżej podpisany Marcin Kurczewski, student Wydziału Matematyki i
Informatyki Uniwersytetu im. Adama Mickiewicza oświadczam, że przedkładaną
pracę dyplomową pt.: "Ataki na kryptograficzne funkcje skrótu" napisałem
samodzielnie.  Oznacza to, że przy pisaniu pracy, poza niezbędnymi
konsultacjami, nie korzystałem z pomocy innych osób, a w szczególności nie
zlecałem opracowania rozprawy lub jej części innym osobom, ani nie odpisywałem
tej rozprawy lub jej części od innych osób.

Oświadczam również, że egzemplarz pracy dyplomowej w formie wydruku
komputerowego jest zgodny z egzemplarzem pracy dyplomowej w formie
elektronicznej.

Jednocześnie przyjmuję do wiadomości, że gdyby powyższe oświadczenie okazało
się nieprawdziwe, decyzja o wydaniu mi dyplomu zostanie cofnięta.


\newpage
\setcounter{tocdepth}{3}
\tableofcontents

\newpage
\pagestyle{fancy}
\fancyhead[L]{\fontsize{10}{0}\selectfont
	Ataki na kryptograficzne funkcje skrótu}
\fancyhead[R]{\fontsize{10}{0}\selectfont
	Marcin Kurczewski}


\section{Wstęp}
Celem tej pracy jest przedstawienie współczesnych metod łamania
kryptograficznych funkcji skrótu, zarówno od strony praktycznej, jak i
teoretycznej. Jako główny temat technologiczny zostały wybrane funkcje
\texttt{MD5} oraz \texttt{SHA-1}.

Praca złożona jest z trzech części. Pierwsza stanowi wprowadzenie do tematyki
kryptograficznych funkcji skrótu, opisując krótko ich rys historyczny,
zgłębiając szczegółowo budowę i przedstawiając ich zastosowania.

Druga część opisuje szereg uniwersalnych technik, które można zastosować do
atakowania większości współczesnych kryptograficznych funkcji skrótu w sposób
niezależny od ich struktury. W części tej są także przedstawione podstawowe
metody obrony przed tego typu atakami.

Część trzecia prezentuje od strony technicznej wybrane wysoce wyspecjalizowane
ataki teoretyczne, jakie powstały na przestrzeni ostatnich lat, pokazując luki
w kryptograficznych funkcjach skrótu.

\newpage
\section{Wprowadzenie do tematyki}
\subsection{Definicja kryptograficznej funkcji skrótu}
Funkcja skrótu jest to funkcja, której dziedziną są ciągi bitów dowolnej
długości, a przeciwdziedziną - ciągi bitów o ograniczonej długości. Rzutowanie
to odbywa się w sposób jednoznaczny, tj. każdemu wejściu funkcja skrótu
przyporządkowuje dokładnie jedno wyjście; formalnie 2zatem jest to surjekcja:

$$ f \colon X \to Y $$
$$ |Y| \leq c, c \in \mathbb{N} $$
$$ \forall{y \in Y} \; \exists{x \in X} \; f(x)=y$$

Kryptograficzna funkcja skrótu to taka funkcja skrótu, która może być
wykorzystywana w zastosowaniach kryptograficznych, a więc wymagających
wysokiego poziomu bezpieczeństwa. Przykładem takiego zastosowania może być np.
szyfrowanie danych.

\subsection{Terminologia użyta w pracy}
Funkcje skrótu często nazywane są "funkcjami haszującymi" niezależnie od
swojej przydatności w dziedzinie kryptologii. W myśl tej tradycji, w
pozostałej części pracy pojęcie "funkcja haszująca" będzie tożsame z
"kryptograficzną funkcją skrótu", gdyż tylko o takich funkcjach ta praca
traktuje.

Dodatkowo, zgodnie z terminologią używaną przez kryptologów, wartość wejściowa
dla funkcji haszującej będzie nazywana "wiadomością", natomiast wartość
wyjściowa funkcji haszującej będzie konsekwentnie nazywana "haszem".

\subsection{Zastosowania funkcji haszujących}
\subsection{Cechy dobrych funkcji haszujących}
\subsection{Rys historyczny}

\section{Konstrukcja funkcji haszujących}
\subsection{Funkcja doskonała}
\subsection{Funkcja jednokierunkowa}
\subsection{Współczesne implementacje}
\subsubsection{Konstrukcja Merkle-Damgarda}
\subsubsection{MD5}
\subsubsection{SHA-1}

\section{Ataki uniwersalne}
\subsection{Atak brutalny}
\subsection{Atak słownikowy}
\subsubsection{Odwrócony atak słownikowy}
\subsection{Tęczowe tablice}
\subsection{Ataki typu Denial of Service}

\section{Ataki teoretyczne}

\section{Podsumowanie}

\end{document}
